\chapter{CONCLUSIONES Y RECOMENDACIONES}

\section{Conclusiones}

\subsubsection{Respuesta a las preguntas de investigación}

Con el fin de consolidar la coherencia científica de la investigación, a continuación se da respuesta sistemática a las preguntas formuladas en el planteamiento del problema (Capítulo 1), fundamentándolas en la evidencia empírica y estadística generada:

\begin{enumerate}
    \item \textit{¿Cómo se manifiestan las sequías en Paraguay en términos de su frecuencia, intensidad, duración y extensión espacial?}
    
    El análisis descriptivo multiescalar demostró que las sequías en Paraguay han experimentado un proceso de agravamiento multidimensional durante el período 2000--2023. En términos de \textit{frecuencia}, se identificó una tendencia positiva, alcanzando un máximo histórico en el año 2020 con 80 meses-bases bajo déficit hídrico moderado o severo. En cuanto a la \textit{intensidad}, los promedios anuales del SPEI-6 se degradaron de valores estables cercanos a $-1.15$ a principios del milenio hacia valores inferiores a $-1.60$, registrándose un mínimo histórico extremo de $-3.7103$ en diciembre de 2023 (base de Mariscal Estigarribia). 
    
    La \textit{duración} promedio de los eventos se incrementó de $2.80$ meses en la primera década a $3.12$ meses en la última (un aumento relativo del 12\% en su persistencia). Finalmente, la \textit{extensión espacial} transitó de un promedio del 10\% de afectación en el año 2000 a proyectarse por encima del 20\% para el año 2023, consolidando a la sequía como un fenómeno de escala nacional.

    \item \textit{¿Existen evidencias estadísticamente significativas de procesos de aridificación en Paraguay, y cómo varía su intensidad entre estaciones climáticas y regiones geográficas?}
    
    Sí, la modelación de datos de panel aportó evidencia del desarrollo de procesos de aridificación en Paraguay, caracterizados por una marcada divergencia estacional y espacial. Mientras que las estaciones de otoño e invierno se mantuvieron hídricamente estables (tendencias de cambio cercanas a cero y estadísticamente no significativas), las épocas cálidas de primavera y verano registraron pendientes negativas hacia la sequedad. 
    
    La mayor tasa de aridificación del país se concentró en la subregión sureste (clúster 1) durante la estación de verano, registrando una tendencia de $-0.04$ puntos de SPEI-6 por año (significativa al 11\% bajo Driscoll-Kraay). Este desplazamiento acumulado de $-0.96$ desviaciones estándar en el balance hídrico estival de la principal zona de producción agropecuaria del país confirma una trayectoria de aridificación estacional localizada.

    \item \textit{¿En qué medida la integración de la temperatura y la precipitación en el análisis del balance hídrico permite detectar condiciones significativamente más deficitarias, atribuibles al aumento de las temperaturas asociadas al calentamiento global?}
    
    La incorporación de la temperatura mediante el índice integrado SPEI-6 reveló que el factor térmico altera sustancialmente la severidad percibida de la sequía en comparación con el análisis tradicional de precipitación univariada (SPI-6). Mientras que las tendencias del SPI-6 resultaron ser no significativas en todas las estaciones (indicando un suministro de lluvia estadísticamente estable), la diferencia de pendientes entre el SPEI-6 y el SPI-6 ($\Delta \beta$) resultó ser significativamente negativa. 
    
    Esto es particularmente evidente en el verano de la subregión centrochaco (clúster 2), donde la diferencia de tendencias ($\Delta \beta = -0.0125$) es altamente significativa al 5\% ($p < 0.05$), y en el verano de la subregión sureste (clúster 1), donde la diferencia ($\Delta \beta = -0.0060$) es significativa al 10\% ($p < 0.10$). Este contraste estadístico prueba que el aumento de las temperaturas y la consecuente demanda evaporativa actúan como el agente motor que intensifica el déficit de agua disponible por encima del régimen de precipitaciones.
\end{enumerate}

\subsubsection{Evaluación y contraste de las hipótesis}

A partir de los resultados empíricos validados de la investigación, se procede a contrastar formalmente las hipótesis de trabajo planteadas en el inicio del estudio:

\begin{itemize}
    \item {Hipótesis 1 (H1): Las sequías en Paraguay durante el período 2000--2023 muestran diferencias significativas en su frecuencia, severidad, duración y extensión de los eventos.}
    
    \textit{Evaluación: Aceptada.} La evidencia analítica-descriptiva del SPEI-6 confirma que todas las dimensiones de la sequía (frecuencia de meses en sequía, intensidad promedio y extrema, persistencia temporal de las rachas de déficit y porcentaje de bases afectadas simultáneamente) han experimentado un incremento sistemático y cuantitativo a lo largo de los 24 años analizados. No obstante, dado que las pendientes de tendencia lineal de largo plazo calculadas mediante el estimador conservador de Driscoll-Kraay mostraron significancia estadística solo para estaciones y regiones específicas, se reconoce que su significancia estadística formal está acotada estacional y espacialmente.

    \item {Hipótesis 2 (H2): El período 2000--2023 evidencia procesos de aridificación estadísticamente significativos en Paraguay con una intensidad variable entre estaciones climáticas y regiones geográficas.}
    
    \textit{Evaluación: Aceptada.} Las tendencias de cambio del balance hídrico estimadas mediante el modelo de datos de panel revelaron una marcada estacionalidad y una divergencia espacial evidente. Mientras que las épocas de otoño e invierno se mantuvieron estables y sin tendencias significativas en ambas regiones, las estaciones cálidas de primavera y verano concentraron las tendencias negativas de aridificación. El punto alto de aridificación se localizó de forma específica en la subregión de mayor relevancia agropecuaria (clúster 1 - sureste) durante la estación de verano, registrando la tasa de cambio negativa más pronunciada del país ($-0.04$ puntos por año, $p < 0.11$).

    \item {Hipótesis 3 (H3): La integración de la temperatura y la precipitación en la evaluación del balance hídrico permite detectar condiciones significativamente más deficitarias que aquellas estimadas únicamente por la precipitación, reflejando el efecto del incremento de las temperaturas asociado al calentamiento global.}
    
    \textit{Evaluación: Aceptada.} La estimación del modelo de diferencias de panel con errores estándar robustos a la correlación espacial demostró que la diferencia de pendientes ($\Delta \beta = \beta_{SPEI} - \beta_{SPI}$) es estadísticamente significativa en el verano de la subregión centrochaco (Clúster 2) al 5\% de significancia ($p < 0.05$), así como en el verano de la subregión sureste (Clúster 1) y en la primavera del Clúster 2 al 10\% de significancia ($p < 0.10$). Dado que la señal pluvial del SPI-6 permaneció estadísticamente estable e insignificante en todos los casos, la significancia de la diferencia negativa confirma que el factor de temperatura y el consecuente incremento de la ETP intensifican de forma significativa el déficit hídrico de las sequías, evidenciando el peso del estrés térmico asociado al cambio climático.
\end{itemize}

\subsubsection{Conclusiones generales}

A partir de la evidencia analítica, climatológica y estadística generada en esta investigación, se desprenden las siguientes conclusiones fundamentales respecto al balance hídrico y las dinámicas de la sequía en Paraguay durante el período 2000--2023:

\begin{enumerate}
    \item La estructuración de una base de datos en formato de panel a partir de los registros diarios de la red oficial de monitoreo terrestre permitió conformar una muestra altamente consistente de 18 bases meteorológicas balanceadas para un período continuo de 24 años. La aplicación de un algoritmo de imputación térmica por capas jerárquicas y la política de no imputación para la variable de precipitación garantizaron la consistencia física de la serie temporal agregada mensual ($N \times T = 5094$ observaciones).
    
    \item El cálculo del índice SPEI a escalas múltiples de acumulación temporal demostró la idoneidad teórica de la distribución Pearson Tipo III para modelar de forma estable el balance hídrico mensual ($D_i$), que incluye valores negativos y ceros. Asimismo, la implementación computacional del recorte de la probabilidad acumulada de la CDF en el entorno $[\epsilon, 1-\epsilon]$ con $\epsilon = 10^{-6}$ eliminó de manera efectiva los valores infinitos en las anomalías climáticas extremas, asegurando la consistencia y el correcto procesamiento de la variable dependiente en los modelos cuantitativos subsiguientes.
    
    \item El análisis analítico-descriptivo confirmó un agravamiento de carácter multidimensional en las características de las sequías en Paraguay. Durante el ciclo 2000--2023, las sequías registraron tendencias positivas en su frecuencia anual (con un pico histórico de 80 meses-bases bajo sequía en el año 2020), una intensificación en el déficit hídrico extremo (mínimo histórico absoluto de $-3.7103$ en diciembre de 2023), una prolongación del 12\% en su duración o persistencia temporal promedio, y una duplicación en la extensión espacial de afectación simultánea a nivel nacional.
    
    \item La regionalización climatológica mediante el agrupamiento jerárquico aglomerativo (HAC) y la optimización del coeficiente de silueta promedio ($s = 0.1758$) determinaron la existencia de dos subregiones climáticas hídricamente homogéneas: la subregión sureste (clúster 1), que comprende la zona de mayor producción agrícola en la región oriental, y la subregión centrochaco (clúster 2), que abarca la región occidental y el nortecentro de la región oriental.
    
    \item La estimación del modelo de regresión de datos de panel mediante el estimador de efectos aleatorios con errores estándar robustos a la correlación espacial de Driscoll-Kraay (1998) aportó evidencia del desarrollo de procesos de aridificación estacional. Se identificó que las estaciones de otoño e invierno se mantienen hídricamente estables en ambas regiones, mientras que las estaciones cálidas de primavera y verano concentran las tendencias negativas de balance hídrico, localizándose la tasa de aridificación estival más intensa en la subregión agrícola del sureste (clúster 1) con una pendiente de $-0.04$ puntos de SPEI-6 por año ($p < 0.11$).
    
    \item El análisis de contraste paralelo entre los modelos de balance hídrico (SPEI-6) y precipitación univariada (SPI-6) demostró la validez del rol dominante del estrés térmico asociado al calentamiento global. Debido a que las lluvias medias (SPI-6) permanecieron estadísticamente estables en todos los casos, la significancia estadística de la diferencia de pendientes ($\Delta \beta < 0$) estimada en el verano del clúster 2 al 5\% de significancia ($\Delta \beta = -0.0125^{**}$) y en el verano del clúster 1 al 10\% ($\Delta \beta = -0.0060^{*}$) comprueba que el agravamiento de las sequías contemporáneas en Paraguay responde al incremento sistemático de la evapotranspiración potencial (demanda atmosférica) inducida por el aumento de las temperaturas máximas extremas, y no a un cambio en el volumen de las precipitaciones.
\end{enumerate}

\section{Recomendaciones}

Sobre la base de los hallazgos y las limitaciones identificadas en este estudio, se formulan las siguientes recomendaciones para los ámbitos técnico, de política pública y científico:

\begin{enumerate}
    \item \textit{Recomendaciones técnicas y de monitoreo climático:} Se sugiere a la Dirección de Meteorología e Hidrología (DMH) de la DINAC fortalecer y modernizar la red de bases meteorológicas terrestres del país. Para superar las limitaciones del método empírico de Thornthwaite en la estimación de la ETP (que depende únicamente de la temperatura media), es indispensable asegurar el registro continuo y de calidad de variables físicas como la velocidad del viento, la radiación solar y los flujos de calor, permitiendo transicionar hacia modelos de balance de agua físicamente más precisos como el estándar de Penman-Monteith (FAO-56) \cite{allen1998crop}.
    
    \item \textit{Recomendaciones para la toma de decisiones y políticas de adaptación:} Se recomienda a los ministerios sectoriales (MAG, MADES) y a los tomadores de decisiones del sector agropecuario integrar métricas sensibles a la temperatura (como el SPEI) en el diseño de planes de contingencia agrícola e infraestructura de riego. Los resultados demuestran que planificar la seguridad hídrica basándose exclusivamente en el monitoreo tradicional de las lluvias (métricas como el SPI) introduce sesgos críticos, ya que un régimen de precipitación estadísticamente estable puede coexistir con un déficit hídrico del suelo severo inducido por olas de calor extremas durante la temporada agrícola.
    
    \item \textit{Futuras líneas de investigación científica:} Se alienta a futuros investigadores a explorar modelos estadísticos no lineales para datos de panel (tales como modelos de panel con umbral) para capturar potenciales quiebres estructurales o cambios bruscos en el régimen del balance hídrico nacional. Asimismo, resulta de gran interés integrar información satelital o de índices de vegetación para correlacionar y validar físicamente la magnitud de los coeficientes de aridificación estimados con los impactos reales de productividad en los cultivos de renta de Paraguay.
\end{enumerate}